\begin{table}[bt]
\centering
\resizebox{\columnwidth}{!}{
\begin{tabular}{|p{5cm}|c|c|c|c|c|c|}
% \hline
\hline % toprule
\multirow{2}{*}{\textbf{Algorithm}} & \multirow{2}{*}{\textbf{\makecell{Syn.\\Sim.}}}
& \multirow{2}{*}{\textbf{\makecell{HDE}}}
& \multicolumn{2}{c|}{\textbf{Predict. Pow.}} 
& \multicolumn{2}{c|}{\textbf{Prob. Solv.}} \\
\cline{4-7}
\textbf{} & \textbf{} &\textbf{} & \textbf{App.} & \textbf{Eff.} & \textbf{Solv.} & \textbf{Fail} \\ % \textbf{Syntactic Similarity} \textbf{Predicted applicability} & \textbf{Predicted effects} \textbf{Plan failed ratio}
\hline  % midrule

% \multicolumn{1}{|c|}{} & \multicolumn{1}{c|}{\textbf{Syntactic}} & \multicolumn{2}{c|}{\textbf{Predictive Power}} & \multicolumn{2}{c|}{\textbf{Problem Solving}}\\
% \hline
% \textbf{Alg + Ref.} & \textbf{Similarity} & \textbf{Applicability} & \textbf{Effects} & \textbf{Solving ratio} & \textbf{Failed ratio} \\ % \textbf{Syntactic Similarity} \textbf{Predicted applicability} & \textbf{Predicted effects} \textbf{Plan failed ratio}
% \hline

% ARMS\footnote{\cite{yang2007learning}} & V & X & X & X & X  \\
ARMS \smallcite{yang2007learning} & \cmark & \xmark & \xmark & \xmark & \xmark & \xmark  \\
\hline
SLAF 
\smallcite{amir2008learning}$^{(*)}$& \xmark & \xmark & \xmark & \xmark & \xmark & \xmark  \\
\hline
FAMA 
\smallcite{aineto2019learning} & \cmark & \xmark & \cmark & \cmark & \xmark & \xmark  \\
\hline
LOCM \smallcite{cresswell2013acquiring} & \xmark & \xmark & (**) & (**) & \xmark & \xmark  \\
\hline
OffLAM \smallcite{LAMANNA2025104256} & \cmark & \xmark & \xmark & \xmark & \cmark & \cmark  \\
\hline
NOLAM \smallcite{Lamanna24} & \cmark & \xmark & \xmark & \xmark & \cmark & \cmark  \\
\hline
% Framer \smallcite{lindsay2017framer} & \cmark\footnotemark[\value{footnote}] & \xmark & \xmark & \xmark & \xmark  \\ RONI: THIS ONE LEARNS FROM NL, NOT FROM SYMBOLIC TRAJECTORIES.
% \hline
SAM \smallcite{juba2021safe} & (***) & \xmark & \xmark & \xmark & \xmark & \xmark  \\
\hline
Cond-SAM \smallcite{mordoch2024safe} & \cmark & \xmark & \xmark & \xmark & \cmark & \xmark  \\ % Conditional-SAM
\hline % \bottomrule
% \textbf{Stochastic SAM~\citep{juba2022learning}} &  &  &  &  &  &  \\ % none of them ? theoretic only ?
% \hline
% \textbf{LATPLAN~\citep{asai2018classical}} &  &  &  &  &  &  \\
% \hline
ROSAME \smallcite{xi2024neuro} & \cmark & \xmark & \xmark & \xmark & \xmark & \xmark   \\
\hline
NL2PDDL\smallcite{Oswald2024DLLMDomainModeling}$^{(****)}$& \cmark & \cmark & \xmark & \xmark & \xmark & \xmark    \\
\hline

\end{tabular}
}
\caption{\small{Evaluation metrics used in prior works on domain model learning. These metrics are defined later in this paper. %Sections~\ref{sec:problem-setting} and~\ref{sec:metrics} \yarin{blank ref}. 
(*) Only reported runtime results. 
(**) Measured how many instances were required to ``converge'', i.e., reached where more data is not helpful, and whether the domain is isomorphic to a reference domain model. (***) Measured the number of trajectories needed until the reference domain model is learned. \roni{Added this "footnote":} (****) They used a Large Language Model (LLM) to learned a domain without knowing the actions' signatures, so they also measured syntactic errors.}}
\label{tab:metric-using-comparison}
\end{table}


% \begin{table}[bt]
% \centering
% \resizebox{\columnwidth}{!}{
% \begin{tabular}{|p{5cm}|c|c|c|c|c|}
% % \hline
% \hline % toprule
% \multirow{2}{*}{\textbf{Algorithm}} & \multirow{2}{*}{\textbf{\makecell{Syn.\\Sim.}}} 
% & \multicolumn{2}{c|}{\textbf{Predict. Pow.}} 
% & \multicolumn{2}{c|}{\textbf{Prob. Solv.}} \\
% \cline{3-6}
% \textbf{} & \textbf{} & \textbf{App.} & \textbf{Eff.} & \textbf{Solv.} & \textbf{Fail} \\ % \textbf{Syntactic Similarity} \textbf{Predicted applicability} & \textbf{Predicted effects} \textbf{Plan failed ratio}
% \hline  % midrule

% % \multicolumn{1}{|c|}{} & \multicolumn{1}{c|}{\textbf{Syntactic}} & \multicolumn{2}{c|}{\textbf{Predictive Power}} & \multicolumn{2}{c|}{\textbf{Problem Solving}}\\
% % \hline
% % \textbf{Alg + Ref.} & \textbf{Similarity} & \textbf{Applicability} & \textbf{Effects} & \textbf{Solving ratio} & \textbf{Failed ratio} \\ % \textbf{Syntactic Similarity} \textbf{Predicted applicability} & \textbf{Predicted effects} \textbf{Plan failed ratio}
% % \hline

% % ARMS\footnote{\cite{yang2007learning}} & V & X & X & X & X  \\
% ARMS \smallcite{yang2007learning} & \cmark & \xmark & \xmark & \xmark & \xmark  \\
% \hline
% SLAF 
% \smallcite{amir2008learning}$^{(*)}$& \xmark & \xmark & \xmark & \xmark & \xmark  \\
% \hline
% FAMA 
% \smallcite{aineto2019learning} & \cmark & \cmark & \cmark & \xmark & \xmark  \\
% \hline
% LOCM \smallcite{cresswell2013acquiring} & \xmark & (**) & (**) & \xmark & \xmark  \\
% \hline
% OffLAM \smallcite{LAMANNA2025104256} & \cmark & \xmark & \xmark & \cmark & \cmark  \\
% \hline
% NOLAM \smallcite{Lamanna24} & \cmark & \xmark & \xmark & \cmark & \cmark  \\
% \hline
% % Framer \smallcite{lindsay2017framer} & \cmark\footnotemark[\value{footnote}] & \xmark & \xmark & \xmark & \xmark  \\ RONI: THIS ONE LEARNS FROM NL, NOT FROM SYMBOLIC TRAJECTORIES.
% % \hline
% SAM \smallcite{juba2021safe} & (***) & \xmark & \xmark & \xmark & \xmark  \\
% \hline
% Cond-SAM \smallcite{mordoch2024safe} & \cmark & \xmark & \xmark & \cmark & \xmark  \\ % Conditional-SAM
% \hline % \bottomrule
% % \textbf{Stochastic SAM~\citep{juba2022learning}} &  &  &  &  &  &  \\ % none of them ? theoretic only ?
% % \hline
% % \textbf{LATPLAN~\citep{asai2018classical}} &  &  &  &  &  &  \\
% % \hline
% ROSAME \smallcite{xi2024neuro} & \cmark & \xmark & \xmark & \xmark & \xmark   \\
% \hline
% NL2PDDL \smallcite{Oswald2024DLLMDomainModeling} & \cmark & \cmark & \cmark & (****) & (****)   \\
% \hline

% \end{tabular}
% }
% \caption{\small{Evaluation metrics used in prior works on domain model learning. These metrics are defined later in this paper. %Sections~\ref{sec:problem-setting} and~\ref{sec:metrics} \yarin{blank ref}. 
% (*) Only reported runtime results. 
% (**) Measured how many instances were required to ``converge'', i.e., reached where more data is not helpful, and whether the domain is isomorphic to a reference domain model. (***) Measured the number of trajectories needed until the reference domain model is learned. (****) They using a different metric }}
% \label{tab:metric-using-comparison}
% \end{table}